Fix \(P\in\mathfrak P_{\mathcal D}^{\mathrm{obs}}\).  Positive unit costs, Assumption~\(\mathrm{ass{:}nondegenerate\text{-}components}\), and Theorem~\(\mathrm{thm{:}budget\text{-}profile}\) give, for every \(d\in\mathcal D\),
\[
 0<v_d
 =\left\{\sqrt{c_{E,d}V_{E,d}}+\sqrt{c_{O,d}V_{O,d}}\right\}^{2}
 <\infty.                                                     \tag{10}
\]
This is a source-score statement on \(\mathfrak P_{\mathcal D}^{\mathrm{obs}}\).  By the gradient interpretation recorded in Definition~\(\mathrm{def{:}bridge\text{-}differentiable\text{-}submodel}\), it is a gradient-variance statement only along the declared bridge-differentiable submodels.

Because \(\mathcal D\) is finite and nonempty, \(\mathcal O_P\) is finite and nonempty; indeed, it contains \((d,d)\) for every \(d\).  By (10), every ratio in (1) is well-defined, positive, and finite.  Thus its maximum exists, is finite, and is at least one.  If the implication on the left of (2) holds, then for every \((d,d')\in\mathcal O_P\), division by \(v_{d'}>0\) yields
\[
 \frac{v_d}{v_{d'}}\le C.
\]
Taking the maximum gives \(C_{\mathrm{ord}}(P)\le C\).  Conversely, if \(C\ge C_{\mathrm{ord}}(P)\) and \(p_d\le p_{d'}\), then \((d,d')\in\mathcal O_P\), so
\[
 v_d\le C_{\mathrm{ord}}(P)v_{d'}\le Cv_{d'}.
\]
This proves (2), including the sharpness of \(C_{\mathrm{ord}}(P)\).

Choose
\[
 d_{\mathrm p}^{\star}\in
 \arg\min_{d\in\mathcal D_{\mathrm{pred}}(P)}v_d,
 \qquad
 d^{\star}\in\arg\min_{d\in\mathcal D}v_d.
\]
Both choices exist because the relevant sets are nonempty subsets of the finite family.  Since \(d_{\mathrm p}^{\star}\) minimizes \(p_d\),
\[
 p_{d_{\mathrm p}^{\star}}\le p_{d^{\star}},
\]
and hence \((d_{\mathrm p}^{\star},d^{\star})\in\mathcal O_P\).  The definitions and (10) give
\[
 1
 \le \frac{v_{d_{\mathrm p}^{\star}}}{v_{d^{\star}}}
 \le C_{\mathrm{ord}}(P).                                   \tag{11}
\]
The middle quantity in (11) is exactly the displayed \(C_{\mathrm{sel}}(P)=\rho_{\mathrm{pred}}(P)\).  Any constant bounding the best-tie rule at this law must be at least this ratio after division by \(v_{d^{\star}}>0\), while this ratio itself gives equality.  It is therefore the least best-tie constant.  Taking suprema over any nonempty \(\mathcal Q\) proves both uniform characterizations: a finite uniform constant exists exactly when the supremum of the corresponding pointwise least constants is finite.  This also proves that the finite preorder maximum in (1), without any additional operator comparison, is the weakest exact lawwise pairwise certificate.  The dependence on \(T_d\) is already carried by the designwise bridges and the resulting values \(v_d\).

Now impose the stated numerical realization of \(\mathsf H_{\mathrm{comp}}\).  Positivity of \(Q_{R,d}\) on \(\mathcal X_d(P)\) makes every quotient in (6) well-defined.  By the definitions of infimum and supremum, (7) is equivalent to the simultaneous inequalities
\[
 \Lambda_-(P)Q_{R,d}(x)
 \le F_d(x)
 \le \Lambda_+(P)Q_{R,d}(x)                                \tag{12}
\]
for every \(d\in\mathcal D\) and \(x\in\mathcal X_d(P)\), with common strictly positive finite endpoints.  Conversely, any common two-sided envelope with constants \(0<L\le U<\infty\) implies \(L\le\Lambda_-(P)\) and \(\Lambda_+(P)\le U\).  Thus no larger common lower constant than \(\Lambda_-(P)\), and no smaller common upper constant than \(\Lambda_+(P)\), can work.  Since every cone is nonempty, its quotient range is nonempty, and hence \(\Lambda_-(P)\le\Lambda_+(P)\).  This proves that (7) is precisely the weakest existence condition for the asserted two-sided envelope and that (6) gives its sharp endpoints.

For any \((d,d')\in\mathcal O_P\), apply (12) at the transported law-specific directions in (4).  The raw-risk ordering gives
\[
\begin{aligned}
 v_d
 &\le \Lambda_+(P)Q_{R,d}\{x_d(P)\}
   =\Lambda_+(P)p_d\\
 &\le \Lambda_+(P)p_{d'}
   =\Lambda_+(P)Q_{R,d'}\{x_{d'}(P)\}\\
 &\le \frac{\Lambda_+(P)}{\Lambda_-(P)}v_{d'}.
\end{aligned}                                                \tag{13}
\]
Maximizing (13) over \(\mathcal O_P\), and then using (11), proves (8).  The exact finite-order argument above did not divide by \(p_d\), so prediction risks may vanish there.  Strict positivity is required only for the conditional quotient representation in (6).

It remains to establish the profiling and generalized-eigenvalue assertions.  Put
\[
 u=c_{E,d}Q_{E,d}(x),
 \qquad
 w=c_{O,d}Q_{O,d}(x).
\]
When \(u,w>0\), direct algebra gives
\[
 \frac{u}{\alpha}+\frac{w}{1-\alpha}
 -(\sqrt u+\sqrt w)^2
 =\frac{\{(1-\alpha)\sqrt u-\alpha\sqrt w\}^{2}}
        {\alpha(1-\alpha)}\ge0.                             \tag{14}
\]
Equality holds at \(\alpha=\sqrt u/(\sqrt u+\sqrt w)\).  If exactly one of \(u,w\) is zero, the same infimum is obtained by taking the appropriate endpoint limit; if both are zero, both sides vanish.  Therefore (14) proves (9) for all nonnegative \(u,w\).

The infimum in (9) need not be a quadratic form.  For an explicit witness, take \(\mathcal H_d=\mathbb R^2\), unit costs, and
\[
 Q_{E,d}(x)=x_1^2,
 \qquad
 Q_{O,d}(x)=x_2^2.
\]
Then \(F_d(x)=(|x_1|+|x_2|)^2\).  With \(e_1=(1,0)\) and \(e_2=(0,1)\),
\[
 F_d(e_1+e_2)+F_d(e_1-e_2)=8,
 \qquad
 2F_d(e_1)+2F_d(e_2)=4,
\]
so the parallelogram identity fails and \(F_d\) is not quadratic.  A largest generalized eigenvalue for a fixed \(\alpha\) can therefore control only the upper quotient associated with that fixed-share quadratic form.  The selection comparison in (13) also needs a strictly positive lower envelope.  Equality further requires transported, law-realized directions that attain the relevant upper and lower envelopes while obeying the raw-risk order.  Comparing design-specific spaces by a single eigenvalue additionally requires a common-space representation and a normalization of the lower envelope.  Observable-law operator bijectivity and raw-risk ordering imply none of those extra properties.  This proves all claims.